\section{\textquestiondown Qué opción utilizar?}

\textbf{Construcción y consultas estáticas.} Los resultados respaldan las
dos primeras hipótesis: Segment Tree se construyó más rápido en todos los
tamaños, mientras que Sparse Table presentó menores tiempos por consulta
en las seis longitudes evaluadas con \(n=10^6\). Para ese tamaño, construir
costó 33,953 ms y 359,880 ms, respectivamente
(tabla~\ref{tab:res-construccion}); en rangos de longitud 511, consultar
costó 507,993 ns y 26,714 ns (tabla~\ref{tab:res-consultas-estaticas}).
Esto es consistente con la preparación \(\Theta(n)\) del árbol frente a
\(\Theta(n\log n)\) de la tabla, y con las consultas RMQ \(\Theta(1)\)
de esta última. No implica que toda consulta del árbol cueste
\(\Theta(\log n)\): los rangos unitarios son constantes en ambas
estructuras. Tampoco demuestra que la tabla amortice su preparación con
cualquier cantidad de consultas.

\textbf{Mantenimiento y carga dinámica.} La tercera hipótesis también
queda respaldada: actualizar Segment Tree fue más barato que reconstruir
Sparse Table en todos los tamaños. Con \(n=10^6\), los tiempos fueron
0,831~\(\mu\)s y 355,673 ms por cambio
(tabla~\ref{tab:res-mantenimiento}). La cuarta hipótesis se cumple para
la carga estudiada: con un cambio y 1.000 consultas por ciclo, el árbol
ganó desde 4.000 elementos; con un millón, necesitó 0,387 ms frente a
355,736 ms. Sin embargo, Sparse Table ganó con 1.000 elementos
(0,105 ms frente a 0,156 ms), por lo que no existe un ganador universal
(tabla~\ref{tab:res-ciclo-dinamico}). El cruce observado entre 1.000 y
4.000 elementos no determina un umbral exacto ni se extrapola a otras cargas.

\textbf{Espacio utilizado.} La quinta hipótesis queda respaldada por
las reservas medidas: con un millón de elementos, Sparse Table retuvo
305,18 MiB frente a 30,52 MiB del árbol, unas diez veces más. Además,
su pico al reconstruir alcanzó 625,61 MiB por la coexistencia de versiones
y reservas temporales; Segment Tree mantuvo 30,52 MiB al actualizar
(tabla~\ref{tab:res-memoria}). Son asignaciones de las estructuras,
no la memoria residente total del proceso.

\textbf{Ventajas y desventajas.} La principal ventaja de Sparse Table
es responder RMQ en tiempo constante y con menores tiempos observados
por consulta; sus desventajas son la preparación y el espacio
\(\Theta(n\log n)\), junto con la ausencia de actualización nativa en
la implementación evaluada, que obliga a reconstruir en este estudio.
Segment Tree ofrece construcción y espacio \(\Theta(n)\), además de
actualizaciones puntuales \(\Theta(\log n)\) sin nuevas reservas.
Su desventaja es un mayor costo de consulta, con peor caso
\(\Theta(\log n)\), que puede acumularse en cargas intensivas de lectura.

\textbf{En resumen.} Conviene Sparse Table para datos inmutables y
memoria suficiente, cuando el ahorro en muchas consultas amortice su preparación.
Conviene Segment Tree si hay cambios frecuentes, memoria limitada o pocas
consultas que no compensen construir la tabla. La decisión debe considerar
construcción, mantenimiento y consultas en conjunto, no solo la operación
más rápida de cada estructura. Estas conclusiones describen las mediciones
y son compatibles con las hipótesis teóricas; no constituyen una prueba
de complejidad asintótica ni de significancia estadística.
